\chapter{Descripción general}
\label{cap:descripcion}

BeeTracer es una plataforma de trazabilidad logística desarrollada por Bilix Ingeniería
(Valparaíso) que digitaliza el ciclo de una operación de desconsolidado en terminales
extraportuarios y bodegas. El comercio exterior chileno opera mayoritariamente con carga
consolidada: la mercancía de varios importadores pequeños viaja en un mismo contenedor, se
declara en un manifiesto electrónico que la nave transmite al Servicio Nacional de Aduanas y,
como el puerto no tiene espacio para abrir contenedores, se separa por cliente en un puerto
seco. Ese momento, el desconsolidado, es donde la custodia de la carga cambia de manos varias
veces en pocas horas y donde históricamente la información se registraba de forma más
precaria: planillas en papel, fotografías dispersas en teléfonos personales y ninguna prueba
del estado de la carga al abrir el contenedor.

La propuesta central de BeeTracer es la \textbf{evidencia fotográfica obligatoria}: la
aplicación no permite avanzar de una etapa del desconsolidado a la siguiente sin capturar el
registro visual correspondiente, de modo que al cerrar la faena existe un expediente completo y
verificable. El sistema tiene una arquitectura cliente-servidor con dos componentes:

\begin{itemize}
    \item \textbf{Plataforma web}, usada desde la oficina para configurar los datos maestros,
    importar manifiestos, programar faenas, validar reportes y consultar el estado de las
    operaciones.
    \item \textbf{Aplicación móvil}, usada en el patio por el jefe tarjador para registrar el
    sello, los lotes separados por cliente, las incidencias y las fotografías. Funciona sin
    conexión y sincroniza la información cuando recupera señal.
\end{itemize}

Desde la perspectiva estructurada, el sistema se entiende como una red de procesos que
transforman flujos de datos. El ciclo completo se descompone en seis procesos principales,
que serán la base de los diagramas de flujo de datos y, más adelante, de los subsistemas del
diagrama de estructura:

\begin{enumerate}
    \item \textbf{Administrar datos maestros}: el Administrador del depósito mantiene los
    usuarios, roles y permisos, tipos de bulto, clientes, trabajadores y cuadrillas.
    \item \textbf{Importar manifiesto}: el Jefe de faena solicita el manifiesto XML de una nave
    y viaje al Servicio Nacional de Aduanas; el sistema lo valida y registra sus contenedores
    y líneas de carga antes del arribo.
    \item \textbf{Programar faenas}: el Jefe de faena planifica el desconsolidado de cada
    contenedor (fecha, franja horaria, cuadrilla y jefe tarjador), y la faena queda disponible
    en la aplicación móvil.
    \item \textbf{Ejecutar operación en terreno}: el Jefe tarjador verifica el sello, registra
    cada lote separado por cliente, las incidencias (daño, faltante, sobrante, mal embalaje,
    sello alterado) y cierra la faena con la fotografía del contenedor vacío. Si no hay
    manifiesto, la carga se ingresa en modalidad \emph{tarja ciega}.
    \item \textbf{Validar reporte}: al cerrar la faena se genera el borrador del reporte; el
    Supervisor del reporte lo revisa y lo aprueba o rechaza la evidencia defectuosa, que se
    recaptura en terreno.
    \item \textbf{Distribuir y consultar reportes}: el reporte aprobado se envía al cliente; el
    Jefe de faena consulta las operaciones y puede reenviar el reporte, y el Cliente del
    depósito consulta el estado de sus operaciones y descarga sus reportes aprobados.
\end{enumerate}

Respecto del Informe 2 se mantienen los mismos actores, reglas de negocio y estados. Solo se
incorporan tres precisiones que el enfoque de datos hace visibles: cada fotografía registra la
\emph{etapa} del desconsolidado a la que pertenece (sello, apertura, lote, incidencia o cierre),
lo que permite ordenar el reporte por etapa y cliente; cada envío del reporte se registra como
un hecho independiente (\emph{envío de reporte}), lo que permite auditar los reenvíos; y,
cuando no existe manifiesto, el Jefe de faena puede registrar manualmente el contenedor al
programarlo, condición necesaria para operar en tarja ciega.
